%% ============================================================
%% 分片 06：原文第 3 节「Topological Field Theory From ${\EUN}=4$ Super Yang-Mills Theory」（原文 glshort.tex 起于第 1216 行）
%% 由 tools/harvmac2tex.py 从 arXiv 源码（hep-th/0604151v3）自动生成。
%% 只译英文散文；公式、编号、\label、\cite 一律不改。
%% ============================================================

\subsection{A Family Of $A$-Models}

\label{sec:familya}

We are in the following general situation.  We have a quantum
theory with a fermionic symmetry $Q=uQ_\ell+vQ_\rr$ that obeys
$Q^2=0$ (modulo a gauge transformation). In such a situation, in
general, the path integral localizes on fields that are invariant
under $Q$. When the space $\eurm Y$ of $Q$-invariant fields is
smooth and the gauge group acts freely on it, the path integral
can be evaluated by a Gaussian or one-loop approximation,
expanding around $\eurm Y$. More generally, one might have to go
to higher order, but the path integral is always determined by the
structure of a finite order infinitesimal neighborhood of $\eurm
Y$.

It suffices to identify the bosonic fields that are invariant
under $Q$, since fermions are in any case infinitesimal, and
automatically are incorporated in perturbation theory.  The
condition for a bosonic field to be invariant under $Q$ is that
$\delta_T\Psi$ vanishes for every fermionic field $\Psi$.  In
other words, $\eurm Y$ is defined by setting to zero the right
hand sides of \eqref{eq:helpme}\ and \eqref{eq:mimo}.

For these expressions to vanish is a condition that is invariant
under scaling of $u$ and $v$ and so only depends on $t=v/u$.
Vanishing of the right hand side of \eqref{eq:helpme}\ gives\footnote{To get the
third equation, we take a $t$-dependent linear combination of  the
conditions $\delta_T\eta=0$ and $\delta_T\tilde\eta=0$, if
$t\not=\pm i$.  What happens if $t=\pm i$ is discussed below.}
% 式 (3.29)
\setcounter{equation}{28}
\begin{equation}\label{eq:loopme}
\begin{aligned}(F-\phi\wedge\phi+tD\phi)^+ & = 0
\\
                      (F-\phi\wedge\phi-t^{-1}D\phi)^-& = 0 \\
                      D^* \phi & = 0.\\\end{aligned}
\end{equation}

 The most elementary interpretation of these equations is
for $-t^{-1}=t$, or equivalently $t=\pm i$.  For these values of
$t$, which will be important for the geometric Langlands program,
let ${\cal A}$ be the ${G}_{\Bbb{C}}$-valued connection ${\cal
A}=A+i\phi$, with curvature ${\cal F}=d{\cal A}+{\cal A}\wedge
{\cal A}$.  Then the first two equations in \eqref{eq:loopme}\ are
equivalent to ${\cal F}=0$, so that a solution of \eqref{eq:loopme}\ at
$t=\pm i$ determines a complex-valued flat connection and hence a
 homomorphism $\pi_1(M)\to G_{{\Bbb{C}}}$. Eqn. \eqref{eq:loopme}\ instructs
us to impose the further condition $D^*\phi=0$ and to divide by
$G$-valued gauge transformations. The combined operation
constructs the moduli space ${\cal Y}$ of homomorphisms
$\vartheta:\pi_1(M)\to {{{G}}}_{\Bbb{C}}$ (we refer to this space
as ${\cal Y}_M$ or ${\cal Y}_M(G)$ if greater clarity is needed).
This statement depends on a theorem of Corlette \cite{ref58} (the two-dimensional case was also
proved by Donaldson in the appendix to \cite{ref20}): dividing by
$G$-valued gauge transformations on the pair $A,\phi$, requiring
that ${\cal F}=0$ and imposing $D^*\phi=0$ is equivalent to simply
dropping that last equation, imposing $\CF=0$ and a condition of
stability\footnote{ As is usual for moduli problems involving
non-compact symmetry groups such as $G_{\Bbb{C}}$, the moduli
space ${\cal Y}_M(G)$ does not literally parametrize all conjugacy
classes of homomorphisms $\pi_1(M)\to G_{\Bbb{C}}$. To get a good
moduli space, one must drop certain ``unstable orbits,''
corresponding in this problem to homomorphisms that can be reduced
to a triangular form but not to a direct sum.  That setting $D^*
\phi=0$ and dividing by $G$-valued gauge transformations has the
effect of dropping such unstable orbits was shown in \cite{ref58}.
We explain the role of stability  more fully in section
第~\ref{sec:complexstr}~节.} and dividing by the group of $G_{\Bbb{C}}$-valued
gauge transformations.

More generally, if $t$ is not real, then the equations \eqref{eq:loopme}\
for fields $A,\phi$ valued in the real Lie algebra of $G$ are
overdetermined, rather than elliptic.  Setting to zero separately
the hermitian and antihermitian parts of these equations, we learn
that they imply ${\cal F}=0$ for any non-real $t$.

We want to understand, however, what happens when $t$ is real.
(The most important values of $t$ for the geometric Langlands
program are $t=\pm i$ and $t=\pm 1$.) In this case, the equations
are real and we get a family of real elliptic equations
parametrized by $t$.  As written, these equations are regular for
$t\not= 0,\infty$.  However, as is clear from the homogeneous
expressions on the right hand side of \eqref{eq:helpme}, the equations can
be extended over $t=0$ or $t=\infty$, by multiplying the second
equation by $t$ or the first by $t^{-1}$.  Thus, at $t=0$, we
should replace the second equation by $(D\phi)^-=0$, and at
$t=\infty$, we replace the first by $(D\phi)^+=0$.  So we get a
family of real elliptic equations parametrized by $\Bbb{RP}^1$.

We can similarly understand the  conditions for unbroken
supersymmetry which come from \eqref{eq:mimo}:
 % 式 (3.30)
\setcounter{equation}{29}
\begin{equation}\label{eq:toopme}
\begin{aligned}D\sigma + t[\phi,\sigma]& = 0\\
                     D\sigma - t^{-1}[\phi,\sigma]& = 0.\\
                     \end{aligned}
\end{equation}

For any $t$ other than $\pm i$, these equations imply that
$D\sigma=[\phi,\sigma]=0$, a condition which means that the gauge
symmetry generated by $\sigma$ leaves invariant the given solution
of \eqref{eq:loopme}.
In addition, for $t\not=\pm i$, a linear combination of the equations $\delta_T\eta=0$
and $\delta_T\tilde \eta=0$ gives
% 式 (3.31)
\setcounter{equation}{30}
\begin{equation}\label{eq:woopem}
[\sigma,\bar\sigma]=0.
\end{equation}

As long as we only consider
supersymmetric fields that are irreducible -- or more generally,
those for which the automorphism group is finite -- \eqref{eq:toopme}\
implies that $\sigma=0$. In this paper, we generally consider only
supersymmetric fields that are irreducible. Some simple remarks on
the  reducible case will appear in \cite{ref26}.

At $t=\pm i$, the conclusion is similar although the details are different.  In this case,
$\eqref{eq:toopme}$ says that $\sigma$ generates a symmetry of the flat $G_\Bbb{C}$-valued connection
$\CA$ or $\bar\CA$.  Again, as long as we only consider supersymmetric fields that are irreducible,
we can assume that $\sigma=0$.  (For $t=\pm i$, in case of a reducible supersymmetric configuration
with $\sigma\not=0$, we do not get separate conditions $D^*\phi=[\sigma,\bar\sigma]=0$, but only
the linear combination $D^*\phi\pm i[\sigma,\bar\sigma]=0$.)

For real $t$, the theory discussed here fits in the general
framework of cohomological field theories \cite{ref59}, like Donaldson theory in four
dimensions or $A$-models in two dimensions.  The fields are the
bosonic fields $A,\phi$ of ${\EUK}=0$; the equations are \eqref{eq:loopme};
and the symmetries are simply the gauge symmetries. As in any such
theory, the path integral in the absence of operator insertions
can be evaluated by counting, with signs, the number of solutions
of the equations, as long as those solutions are isolated and
irreducible.  The sign with which a given solution contributes is
simply  the sign of the corresponding one-loop fermion
determinant. A smooth compact family of irreducible solutions,
parametrized by a space ${\eurm Y}_0$ (which is a component of
${\eurm Y}$), makes a contribution that is plus or minus the Euler
characteristic of ${\eurm Y}_0$, with the sign again coming from
the fermion determinant. Contributions of singularities in ${\eurm
Y}$ and of reducible solutions -- such as the trivial solution
with $A=\phi=0$ -- are more subtle to evaluate and will not be
considered here.

Since the number of solutions, weighted by sign, is invariant
under continuous deformation of a family of elliptic equations,
the partition function on a closed four-manifold without operator
insertions must be independent of $t$, at least for real $t$. By
holomorphy, this is also true for complex $t$.  That the partition
function on a closed four-manifold is independent of $t$ can be
seen more directly from the path integral, as we explain later.
{}From a mathematical point of view, the vanishing theorems of
section 第~\ref{sec:vanishing}~节\ will give a strong statement about
$t$-independence: the space of solutions of \eqref{eq:loopme}\ is actually
independent of $t$ except at $t=0,\infty$.

The importance of the parameter $t$ is that it is possible to
introduce operators or, in case $M$ has a boundary, boundary
conditions, that preserve the topological symmetry only for a
particular value of $t$. Like Donaldson theory, the TQFT
considered here has local operators that preserve the topological
symmetry. They have been discussed in \cite{ref53,ref54,ref55,ref56,ref57} and
will be further analyzed elsewhere \cite{ref26}.  However, they are not
the most important operators for the  geometric Langlands program.
The important operators, introduced in section 第~\ref{sec:loopop}~节, will be
the Wilson and 't Hooft line operators, which have no obvious
close analogs in more familiar cohomological field theories.

If our theory for real $t$ is somewhat analogous to a
two-dimensional $A$-model, what does it correspond to for complex
$t$?  The theory at $t=\pm i$ is analogous to a two-dimensional
$B$-model.
 In the $B$-model that describes  maps $\Phi:\Sigma\to X$ from a
Riemann surface $\Sigma$ to a complex manifold $X$,  the
supersymmetric fields are the constant maps of $\Sigma$ to $X$.
The obvious similarity between the two cases is that the constant
maps or the flat $G_{\Bbb{C}}$-valued connections are relatively
elementary to describe  by using over-determined rather than
elliptic equations (namely $d\Phi=0$ in one case, ${\cal F}=0$ in
the other) in which the complex structure is obvious. A deeper
analogy will be apparent when we consider dimensional reduction in
section 第~\ref{sec:compacthitchin}~节.

\subsection{Vanishing Theorems}

\label{sec:vanishing}

Let $\EUV^+(t)= (F-\phi\wedge\phi+tD\phi)^+$, $\EUV^-(t)=
(F-\phi\wedge\phi-t^{-1}D\phi)^-$, and $\EUV_0= D^* \phi$, so the
supersymmetric equations are $\EUV^+(t)=\EUV^-(t)=\EUV_0=0$.  In
this section, we concentrate on the case of real $t$, so that the
supersymmetric equations are elliptic.

Like many first order equations associated with supersymmetry,
such as the equations for a holomorphic curve, the Yang-Mills
instanton equations, the Seiberg-Witten equations, or the
equations studied in \cite{ref45}, these are subject to unusual
vanishing theorems. In the case at hand, we have:

\par\medskip
 {\it Vanishing Theorem 1~}  Let $M$ be a compact
four-manifold  without boundary, and $E$ a $G$-bundle over $M$
with
 nonzero Pontryagin class, $\int_M
\Tr\,F\wedge F\not=0$.  Then for any $t\not=0,\infty$, the
supersymmetric equations $\EUV^+(t)=\EUV^-(t)=\EUV_0=0$ have no
solutions.

\par\medskip  {\it Vanishing Theorem 2~} On such an $M$, any
field that obeys those equations
 for one value of $t$ other than $0,\infty$ obeys them for all $t$
 and hence is given by a flat $G_\Bbb{C}$-valued connection.

\par\medskip

The shortest proof of the first vanishing theorem comes from the
following identity: % 式 (3.32)
\setcounter{equation}{31}
\begin{equation}\label{eq:hombo}
\int_M\Tr\,F\wedge
F=\int_M\Tr\left(\EUV^+(t)\wedge \EUV^+(-t^{-1})+\EUV^-(t)\wedge
\EUV^-(-t^{-1})\right).
\end{equation}
 The proof of the identity is
straightforward and uses integration by parts. Clearly, this
identity implies that if for some $t$, $\EUV^\pm(t)=0$, then
$\int_M\Tr\, F\wedge F=0$.

This proof fails if $t=0,\infty$ because of the poles in
$\EUV^\pm$. In fact, the result is false for $t=0,\infty$; the
equations at $t=0$ or $t=\infty$ have at least the obvious
solutions given by an instanton or anti-instanton with $\phi=0$.

For the second vanishing theorem, we need a similar but more
intricate identity: % 式 (3.33)
\setcounter{equation}{32}
\begin{equation}\label{eq:ormob}
\begin{aligned}&-\int_Md^4x\sqrt
g\,\Tr\,\left({t^{-1}\over
t+t^{-1}}\EUV^+(t)_{\mu\nu}\EUV^+(t)^{\mu\nu}+{t\over
t+t^{-1}}\EUV^-(t)_{\mu\nu}\EUV^-(t)^{\mu\nu}+\EUV_0^2\right)\\
&=-\int d^4x\sqrt g\,\Tr\,\left({1\over 2}F_{\mu\nu}
F^{\mu\nu}+D_\mu\phi_\nu D^\mu\phi^\nu
+R_{\mu\nu}\phi^\mu\phi^\nu+ {1\over
2}[\phi_\mu,\phi_\nu]^2\right)\\&~~~~~~~+{t-t^{-1}\over
(t+t^{-1})}\int_M\Tr\, F\wedge F. \\\end{aligned}
\end{equation}
 What is surprising about
this identity is that apart from the topological invariant
proportional to $\int\,\Tr\,F\wedge F$, the right hand side is
independent of $t$.  The proof is similar to the proof of \eqref{eq:hombo},
but more elaborate. $R_{\mu\nu}$ is the Ricci tensor of $M$, which
enters in integrating by parts to relate $\int_Md^4x\sqrt g
D_\mu\phi_\nu D^\nu\phi^\mu$ to $\int d^4x\sqrt g
(D_\mu\phi^\mu)^2$.  Unlike \eqref{eq:hombo}, \eqref{eq:ormob}\ has a natural limit as
$t\to 0,\infty$.  One simply replaces $t (\EUV^-)^2/(t+t^{-1})$ by
$\lim_{t\to 0}(t\EUV^-)^2$ or $t^{-1}(\EUV^+)^2/(t+t^{-1})$ by
$\lim_{t\to \infty}(t^{-1}\EUV^+)^2$.  So arguments based on
\eqref{eq:ormob}\ are valid at $t=0,\infty$.

In section 第~\ref{sec:toplag}~节, we will understand the physical meaning of
\eqref{eq:ormob}. For now, we simply use it to complete the proof of the
vanishing theorems.

To prove the second vanishing theorem, we must show that if the
bundle $E$ admits a pair $A,\phi$ for which $\EUV^\pm(t)=\EUV_0=0$
for some value of $t$ other than $0,\infty$, then these equations
hold for all $t$.  We already know that in such a solution,
$\int_M\Tr\,F\wedge F=0$, whence the right hand side of \eqref{eq:ormob}\ is
independent of $t$.  So if the left hand side vanishes for one
value of $t$, it vanishes for all $t$. Moreover, the left hand
side is a sum of squares, so it vanishes for given $t$ if and only
if $\EUV^\pm(t)=\EUV_0=0$. Moreover, as we have already noted, the
vanishing of $\EUV^\pm(t)$ and $\EUV_0$ for all $t$ is equivalent
to the vanishing of the curvature ${\cal F}=d{\cal A}+{\cal
A}\wedge {\cal A}$ of the $G_{\Bbb{C}}$-valued connection ${\cal
A}=A+i\phi$.  This completes the argument.

Finally, we can use \eqref{eq:ormob}\ to prove a slightly stronger version
of the first vanishing theorem:

\par\medskip{\it Vanishing Theorem 1$'$~} For $\int_M\Tr\,
F\wedge F>0$, the supersymmetric equations have no solutions
except at $t=0$, and for $\int_M\Tr\,F\wedge F<0$, they have no
solutions except at $t=\infty$.\par\medskip

Let ${\eurm B}(t)$ denote the left hand side of \eqref{eq:ormob}, and let
$f(t)=(t-t^{-1})/(t+t^{-1})$.  We have for any $t,u$
% 式 (3.34)
\setcounter{equation}{33}
\begin{equation}\label{eq:toto}
{\eurm B}(u)-{\eurm B}(t)=(f(u)-f(t))\int_M\Tr \,F\wedge
F.
\end{equation}
 If $\EUV^\pm(t)=\EUV_0=0$, then ${\eurm B}(t)=0$, and as the
function ${\eurm B}$ is positive semidefinite, the left hand side
of \eqref{eq:toto}\ is positive semidefinite. To make the right hand side
positive semidefinite, $t$ must be a minimum of the function $f$
if $\int_M\Tr\, F\wedge F>0$, and a maximum if $\int_M\Tr\,F\wedge
F<0$.  But the only minimum is at $t=0$, and the only maximum is
at $t=\infty$.

\par\medskip{\it A Note On Compactification}

A special case of the vanishing theorems  is relevant to
compactification.  We take $M=\Sigma\times C$, where $\Sigma$ and
$C$ are two compact Riemann surfaces.  In our applications, $C$ is
the Riemann surface on which we study the geometric Langlands
program, and we think of $\Sigma$ as being much larger than $C$.
However, in topological field theory, the metrics on $\Sigma$ and
$C$ are inessential.

For $t\not=0,\infty$, the second vanishing theorem says that any
solution of the supersymmetric equations on $M$ is given by a flat
$G_{\Bbb{C}}$-valued connection.  A flat bundle $E\to
M=\Sigma\times C$ has a decomposition\footnote{In this discussion, for
simplicity, we use the language of vector bundles, as if
$G=SU(N)$. To extend the argument to any $G$, one simply
formulates it in terms of subgroups of $G$ rather than subbundles
of $E$.} $E=\oplus_{i=1}^n E'_i\otimes E''_i$, where the $E'_i$
are flat bundles pulled back from  $\Sigma$ and the $E''_i$ are
flat bundles pulled back from $C$.

If we restrict to a point $x\in \Sigma$, then the $E'_i$ become
trivial. If $d_i$ is the rank of $E_i'$, then $E$ reduces to
$E_x=\oplus_{i=1}^n d_iE''_i$, a direct sum of flat bundles pulled
back from $C$. $E_x$ is a flat bundle over $C$, and as such
defines a point in Hitchin's moduli space $\MH$, which we describe
more thoroughly in section 第~\ref{sec:compacthitchin}~节. For now, we can just think of $\MH$ as the
moduli space of flat $G_\Bbb{C}$-bundles over $C$.  Thus any
supersymmetric field on $\Sigma\times C$ is given by a map from
$\Sigma$ to $\MH$, namely $x\to E_x$.

The moduli space $\MH$ has singularities that correspond to
reducible flat bundles.  In the present construction, the bundle
$E_x$ is reducible if $n>1$. For $n=1$, we have $E=E'\otimes E''$,
where the two factors are pullbacks from $\Sigma $ and $C$,
respectively.  $E_x$ is still reducible unless $E'$ has rank 1, in
which case, for $G=SU(N)$, it must be trivial.  Thus, the case in
which the singularities of $\MH$ are avoided is precisely the case
that $E$ is the pullback of an irreducible flat bundle $E''\to C$.

 The supersymmetric fields are
also the configurations that minimize the action; this will become
clear in section 第~\ref{sec:toplag}~节, when we construct the supersymmetric
Lagrangian.   In the above reasoning, we took $\Sigma$ and $C$ to
be compact.  In our applications, we are also interested in the
case that $\Sigma$ is a complete but noncompact two-manifold, such
as $\Bbb{R}^2$, $\Bbb{R}^2_+$, or $\Bbb{R}\times I$ ($\Bbb{R}^2_+$
is a halfspace and $I$ is a closed interval). Given appropriate
boundary conditions and asymptotic behavior at infinity to justify
the proof of the second vanishing theorem, one still has in these
cases the decomposition of supersymmetric fields
$E=\oplus_{i=1}^nE'_i\otimes E''_i$. (For $\Sigma$ flat and
simply-connected, by $E'_i$ we mean a bundle with trivial
connection and covariantly constant Higgs field.) This leads to
the same conclusion that we had for compact $\Sigma$: a
supersymmetric field that avoids singularities of $\MH$ is a
pullback from $C$.
